\chapter{Manifold-Guided Adversarial Engine (\mgae{})}
\label{ch:mgae}

\section{Motivation: The Manifold Problem}

Standard adversarial perturbations (FGSM, PGD) perturb features in the
gradient direction without constraint. For network traffic, this
frequently produces samples outside the support of any real distribution:
negative packet counts, byte rates exceeding link capacity, zero-duration
flows with non-zero bytes. Such samples are trivially flagged by
distribution-aware anomaly detectors even when they fool classifier-based IDS.

\mgae{} addresses this by learning the benign traffic manifold
$\mathcal{M}_{\mathrm{benign}}$ via a denoising diffusion process
and projecting attack samples onto $\mathcal{M}_{\mathrm{benign}}$
before submission to the target IDS.

\section{Cosine Diffusion Schedule}

Following \cite{nichol2021}, we use a cosine noise schedule:
\begin{equation}
  \bar\alpha_t = \frac{\cos^2\!\left(\frac{t/T + s}{1+s}\cdot\frac{\pi}{2}\right)}
    {\cos^2\!\left(\frac{s}{1+s}\cdot\frac{\pi}{2}\right)},
  \quad s = 0.008,\quad T = 30
\end{equation}
This schedule avoids the abrupt noise increase at $t=0$ seen in linear
schedules, which is particularly important for low-dimensional tabular data.

\section{Score Network}

The score network $\boldsymbol\epsilon_\theta(\mathbf{x}_t, t)$ predicts
the noise $\boldsymbol\epsilon$ added at step $t$:
\begin{equation}
  \mathcal{L}_{\mathrm{score}} = \E_{t, \mathbf{x}_0, \boldsymbol\epsilon}
    \left[\|\boldsymbol\epsilon - \boldsymbol\epsilon_\theta(\mathbf{x}_t, t)\|^2\right]
\end{equation}
Architecture: input $[\mathbf{x}_t \in \mathbb{R}^d;\, \mathbf{e}_t \in \mathbb{R}^8]$
where $\mathbf{e}_t$ is a sinusoidal time embedding;
three hidden layers of width 64 with ReLU; output in $\mathbb{R}^d$.

Training: 25--40 epochs on benign traffic samples, using an evolutionary
strategy (ES) gradient-free optimiser due to the numpy-only implementation.
In the PyTorch version, standard Adam optimisation is used.

\section{MGAE Perturbation Procedure}

\begin{algorithm}[H]
\caption{\mgae{} Perturbation}
\label{alg:mgae}
\begin{algorithmic}[1]
\Require Attack sample $\mathbf{x}_{\mathrm{atk}}$, trained score net
  $\boldsymbol\epsilon_\theta$, injection step $t_{\mathrm{inj}}$,
  interpolation weight $\alpha$
\State Normalise: $\hat{\mathbf{x}} = \mathrm{clip}
  \!\left(\frac{\mathbf{x}_{\mathrm{atk}} - \boldsymbol\mu}
  {3\boldsymbol\sigma}, -1, 1\right)$
\State Add noise: $\mathbf{x}_{t_{\mathrm{inj}}} \sim
  q(\mathbf{x}_{t_{\mathrm{inj}}} \mid \hat{\mathbf{x}})$
\For{$t = t_{\mathrm{inj}}, \ldots, 1$}
  \State $\hat{\mathbf{x}}_0 = \mathrm{clip}\!\left(
    \frac{\mathbf{x}_t - \sqrt{1-\bar\alpha_t}\,\boldsymbol\epsilon_\theta(\mathbf{x}_t,t)}
    {\sqrt{\bar\alpha_t}}, -1, 1\right)$
  \State $\mathbf{x}_{t-1} = \sqrt{\bar\alpha_{t-1}}\,\hat{\mathbf{x}}_0
    + \sqrt{1-\bar\alpha_{t-1}}\,\boldsymbol\epsilon_\theta(\mathbf{x}_t,t)$
\EndFor
\State Interpolate: $\mathbf{x}_{\mathrm{adv}} = \alpha\,\mathbf{x}_0
    + (1-\alpha)\,\hat{\mathbf{x}}$
\State Preserve high-intensity features:
  $\mathbf{x}_{\mathrm{adv}}[\hat{\mathbf{x}}>0.5]
    \leftarrow 0.55\,\hat{\mathbf{x}} + 0.45\,\mathbf{x}_{\mathrm{adv}}$
\State De-normalise and return $\mathbf{x}_{\mathrm{adv}}$
\end{algorithmic}
\end{algorithm}

\section{Evasion--Distortion Trade-off}

The injection step $t_{\mathrm{inj}}$ controls the evasion--distortion
trade-off: higher $t_{\mathrm{inj}}$ adds more noise and produces
stronger manifold alignment (better evasion) at the cost of larger
perturbation magnitude. \Cref{fig:mgae_sensitivity} (\cref{ch:experiments})
shows that $t_{\mathrm{inj}} = 10$ achieves the optimal trade-off,
maximising evasion probability while keeping $\|\boldsymbol\delta\|_2 < 6.0$.
